\section*{अध्याय \ref{ch:b3:rna-regulation} --- अकूटलेखी RNA और अनुलेखनोत्तर नियमन}

\begin{solution}{exo:b3:rna-regulation:1}
miRNA: जीनोम में हेयरपिन के रूप में कूटबद्ध, \qty{22}{nt}, पहले \omterm{def:b3:rna-regulation:mirna}{ड्रोशा}
फिर \omterm{def:b3:rna-regulation:rnai}{डाइसर} से कटा, अपने बीज से अपूर्ण युग्मन करता है और अनेक
संदेशवाहकों के अनुवादन तथा स्थायित्व को दमित करता है। \omterm{def:b3:rna-regulation:ncrna}{siRNA}: लंबे
द्विलड़ी RNA से (विषाणुजन्य, पारजीनी, प्रायोगिक), \qty{21}{nt},
डाइसर-निर्भर, पूर्ण युग्मन करता है और \omterm{def:b3:rna-regulation:rnai}{आर्गोनॉट} को लक्ष्य विदलित करने
का निर्देश देता है। \omterm{def:b3:rna-regulation:ncrna}{piRNA}: \omterm{def:b3:rna-regulation:pirna}{piRNA गुच्छों} के एकलड़ी अनुलेखों से,
\qtyrange{26}{31}{nt}, डाइसर-निरपेक्ष, जनन-वंश में PIWI प्रोटीनों पर
भरा जाता है, ट्रांसपोज़ॉन अनुलेखों को विदलित करता है (पिंग-पॉन्ग) और
ट्रांसपोज़ॉन DNA के \omterm{def:b3:chromatin-epigenetics:nucleosome}{क्रोमैटिन} मौनीकरण का निर्देश देता है।
\end{solution}

\begin{solution}{exo:b3:rna-regulation:2}
Pol II pri-miRNA का अनुलेखन करती है; \omterm{def:b3:rna-regulation:mirna}{ड्रोशा} (DGCR8 के साथ) केंद्रक में
हेयरपिन काटकर निकालती है; एक्सपोर्टिन-5 pre-miRNA का निर्यात करता है;
\omterm{def:b3:rna-regulation:rnai}{डाइसर} फंदा काट देती है और \qty{22}{nt} का युग्मक बचता है; एक लड़ी
\omterm{def:b3:rna-regulation:rnai}{आर्गोनॉट} में भरी जाती है; बीज (न्यूक्लियोटाइड 2--8) किसी 3$'$ UTR स्थल
से युग्मित होता है; \omterm{def:b3:rna-regulation:rnai}{आर्गोनॉट} अपऐडीनिलेज़ तथा विटोपन तंत्र बुलाता है और
आरंभन रोकता है --- संदेशवाहक कम अनुवादित होता है और तेज़ी से क्षय होता है।
\end{solution}

\begin{solution}{exo:b3:rna-regulation:3}
किसी फ़ेज पॉलीमरेज़ से पात्रे बने RNA में द्विलड़ी दूषक थे (पॉलीमरेज़
साँचे की दूसरी लड़ी की भी नकल करती है और सिरों से आगे निकल जाती है),
इसलिए ``संवेदी'' तैयारियाँ असली उत्प्रेरक ढो रही थीं। फ़ायर और मेलो ने
हर लड़ी शुद्ध की, दिखाया कि अकेली किसी से बहुत कम होता है, फिर उन्हें
जानबूझकर तापानुशीतित किया: द्विलड़ी RNA कम-से-कम सौ गुना अधिक प्रबल था,
और प्रति कोशिका कुछ ही अणु पर्याप्त थे।
\end{solution}

\begin{solution}{exo:b3:rna-regulation:4}
लौह-भुखमरी में: IRP, IRE से बँधता है। फ़ेरिटिन: अनुवादन का आरंभन रुक
जाता है, कोई भंडारण-प्रोटीन नहीं बनता (अनुवादन का नियंत्रण)।
ट्रांसफ़ेरिन ग्राही: संदेशवाहक अपनी न्यूक्लिएज़ से सुरक्षित होकर जमा
होता है, इसलिए अधिक ग्राही बनते हैं और अधिक लोहा आयात होता है
(संदेशवाहक के स्थायित्व का नियंत्रण)।
\end{solution}

\begin{solution}{exo:b3:rna-regulation:5}
$\gamma = \ln 2/\qty{120}{min} = \qty{0.0058}{min^{-1}}$; $m^{*} =
5/0.0058 \approx 870$ अणु। क्षय दोगुना: $m^{*} \approx 430$,
अर्ध-काल $\qty{120}{min}/2 = \qty{60}{min}$।
\end{solution}

\begin{solution}{exo:b3:rna-regulation:6}
कुल 3$'$ UTR अनुक्रम $2\times 10^{7}$ nt; संयोगी मेल $2\times
10^{7}/4^{7} = 2\times 10^{7}/\num{16384} \approx \num{1200}$ स्थल।
किसी वास्तविक लक्ष्य की पहचान जातियों के पार स्थल के संरक्षण से होती
है, अनुकूल संदर्भ से (AU-समृद्ध पड़ोस, विराम-कूटक से दूर का स्थल,
miRNA के 3$'$ सिरे पर अतिरिक्त युग्मन), miRNA तथा लक्ष्य के एक ही
कोशिका में सह-अभिव्यक्त होने से, और प्रयोग से: miRNA के होने पर
संदेशवाहक गिरता है और स्थल उत्परिवर्तित होने पर नहीं गिरता।
\end{solution}

\begin{solution}{exo:b3:rna-regulation:7}
(क) XX, \emph{tra} अकारक: SXL बनता है पर TRA नहीं, \emph{dsx}
चूक-विच्छेदन यानी नर विच्छेदन लेता है --- कायिक रूप से नर (बंध्य)।
(ख) अनवरत TRA वाला XY: TRA-2 दोनों लिंगों में है, इसलिए DSX मादा रूप
में बनता है --- कायिक रूप से मादा परिवर्धन (एक बंध्य
``छद्म-मादा'', जिसका जनन-वंश दूसरे संकेतों का अनुसरण करता है)। (ग)
\emph{dsx} नहीं: DSX का कोई भी रूप नहीं, इसलिए अंतिम विभेदन का कोई भी
कार्यक्रम थोपा नहीं जाता --- मिश्रित या अधूरे लक्षणों वाला अंतर्लिंगी।
\end{solution}

\begin{solution}{exo:b3:rna-regulation:8}
स्थायी अवस्था पर संदेशवाहक, और इसलिए प्रोटीन, संदेशवाहक के अर्ध-आयु
काल के समानुपाती है: $\qty{180}{min}/\qty{15}{min} = 12$ गुना अधिक
प्रोटीन। प्रोटीन एक सामान्य वृद्धि-संकेत है जिसे क्षणिक रहना चाहिए,
हर उद्दीपन के बाद एक स्फुरण; बारह गुना और लगातार बनने पर वह बिना किसी
उद्दीपन के प्रचुरोद्भवन चलाता है। उसे कर्कटजनक मात्रा और समय बनाते
हैं, अनुक्रम नहीं।
\end{solution}

\begin{solution}{exo:b3:rna-regulation:9}
(क) प्रयोगशाला माता, वन्य पिता: अंडों में P अवयवों के विरुद्ध कोई
\omterm{def:b3:rna-regulation:ncrna}{piRNA} नहीं (माता के वंशक्रम ने उन्हें कभी नहीं देखा), पैतृक P अवयव
संतति के जनन-वंश में स्वच्छंद स्थानांतरित होते हैं --- बंध्य
(दुर्जनन)। (ख) वन्य माता: उसके अंडों में P-अवयव \omterm{def:b3:rna-regulation:ncrna}{piRNA} हैं, जो संतति
में अवयवों को मौन करते हैं --- उर्वर। विषमता \omterm{def:b3:rna-regulation:ncrna}{piRNA} के मातृ जमाव से है,
यानी RNA की वंशागति से, जीनों की नहीं।
\end{solution}

\begin{solution}{exo:b3:rna-regulation:10}
प्रोटीन अपने ही जीवन-काल $\tau_{p}$ पर अपने संदेशवाहक का समाकलन करता है।
संदेशवाहक $\tau_{m} = 1/(\gamma + \gamma' R)$ जीते हैं; $\tau_{p}$ के
दौरान प्रोटीन इसलिए लगभग $\tau_{p}/\tau_{m}$ स्वतंत्र संदेशवाहक
जीवन-काल देखता है, जिनमें से हर एक कोई यादृच्छिक जन्म--मृत्यु घटना है।
$N$ स्वतंत्र घटनाओं पर लिए गए औसत का सापेक्ष उतार-चढ़ाव
$1/\sqrt{N}$ की तरह गिरता है, इसलिए \omterm{def:b3:rna-regulation:ncrna}{माइक्रो RNA} के नियंत्रण में छोटा
$\tau_{m}$ --- उसी माध्य संदेशवाहक पर, जिसके लिए समानुपाती रूप से ऊँचा
$k$ चाहिए --- अधिक चिकना प्रोटीन देता है। कोशिका \omterm{def:b3:rna-regulation:ncrna}{माइक्रो RNA} और
अतिरिक्त अनुलेखन का मूल्य परिशुद्धि और गति खरीदने के लिए चुकाती है:
तेज़ संश्लेषण और तेज़ क्षय के संतुलन से तय स्तर उसी स्तर की तुलना में
कम कोलाहली और जल्दी पुनःसमायोजित होने वाला है जो धीमे संश्लेषण और धीमे
क्षय से तय हो।
\end{solution}

\begin{solution}{exo:b3:rna-regulation:11}
कृमि: RNA-निर्भर RNA पॉलीमरेज़ स्वयं लक्ष्य संदेशवाहक से नए \omterm{def:b3:rna-regulation:ncrna}{siRNA}
बनाती है, इसलिए जब तक लक्ष्य अनुलेखित होता रहे संकेत नया होता रहता है,
विभाजनों के तनुकरण से बचता है और, कोशिकाओं के बीच द्विलड़ी-RNA वाहिका
के साथ, दूसरे ऊतकों में तथा कुछ पीढ़ियों के लिए जनन-वंश में फैलता है।
स्तनधारी: कोई प्रवर्धन नहीं, इसलिए \omterm{def:b3:rna-regulation:ncrna}{siRNA} हर विभाजन पर खपते और तनु होते
हैं, केवल उन्हीं कोशिकाओं में काम करते हैं जिन्हें वे मिले, और
विभाजनशील कोशिकाओं में दिनों में मिट जाते हैं। परिणाम: किसी \omterm{def:b3:rna-regulation:ncrna}{siRNA} औषध
को हर लक्ष्य कोशिका में पहुँचाना पड़ता है, वह उन कोशिकाओं में सबसे
अच्छा काम करती है जो विभाजित नहीं होतीं (यकृत-कोशिकाएँ, इसीलिए यकृत की
औषधें), और उसे बार-बार देना पड़ता है।
\end{solution}

\begin{solution}{exo:b3:rna-regulation:12}
(क) प्रकार्यात्मक RNA: अनुलेख बन जाने के बाद उसे नष्ट करने से (\omterm{def:b3:rna-regulation:ncrna}{siRNA},
प्रतिसंवेदी ऑलिगोन्यूक्लियोटाइड) कोई लक्षण-प्ररूप मिलता है, और जीनोम
में कहीं और से किसी पारजीन द्वारा RNA अभिव्यक्त करने पर वह ठीक हो जाता
है। (ख) प्रकार्यात्मक अनुलेखन: कोई अगेती बहुऐडीनिलन संकेत डालकर
अनुलेखन रोकने से (या प्रवर्तक हटाने से) लक्षण-प्ररूप मिलता है, जबकि
RNA को अनुलेखनोत्तर नष्ट करने से नहीं मिलता, और कोई बाह्यस्थानी पारजीन
उसे ठीक नहीं करता --- अनुलेखन की क्रिया (या DNA अवयव) मायने रखती है,
जैसा अनेक संवर्धक-सहचारी RNA के लिए है। (ग) कोलाहल: कोई भी हस्तक्षेप
कुछ नहीं बदलता, RNA प्रति कोशिका एक प्रति से कम है, उसका अनुक्रम और
अभिव्यक्ति संरक्षित नहीं हैं।
\end{solution}

\begin{solution}{pb:b3:rna-regulation:1}
\textbf{1.} $\gamma = \ln 2/40 = \qty{0.0173}{min^{-1}}$; $m^{*} =
3/0.0173 \approx 170$ संदेशवाहक।
\textbf{2.} $\delta_{p} = \ln 2/180 = \qty{0.00385}{min^{-1}}$; $P^{*} =
\beta m^{*}/\delta_{p} = 2\times 173/0.00385 \approx 9.0\times 10^{4}$
प्रोटीन।
\textbf{3.} क्षय $4\gamma = \qty{0.069}{min^{-1}}$; $m^{*} = 3/0.069
\approx 43$; काल-नियतांक $1/0.069 = \qty{14}{min}$।
\textbf{4.} $\beta = 1$: $P^{*} = 43/0.00385 \approx 1.1\times 10^{4}$;
दमन $8$ गुना ($4$ क्षय से, $2$ अनुवादन से)।
\textbf{5.} संदेशवाहक का अर्ध-काल $\ln 2/(4\gamma) = \qty{10}{min}$।
प्रोटीन फिर अपने ही अर्ध-आयु काल \qty{3}{h} से क्षय होता है: स्विच को
प्रोटीन का अपघटन सीमित करता है।
\textbf{6.} हाँ: संदेशवाहक केवल $1/4$ तक गिरता है (अब भी सहज ही
पकड़ में आता है) जबकि प्रोटीन $1/8$ तक गिरता है और पुराना प्रोटीन
बदलते-बदलते और गिरता जाता है; यह आरंभिक प्रेक्षण कि प्रोटीन ``जब mRNA
बना हुआ था तब'' लुप्त हो गया, अधिकतर दमन के अनुवादनी आधे भाग को
दर्शाता था।
\textbf{7.} $10^{6}/250 = 4000$ अणु प्रति अंडा।
\textbf{8.} अंडे से निकलते समय $4000/550 \approx 7$ प्रति कोशिका।
$4000/2^{n} < 1$ के लिए युग्मनज से $n \ge 12$ विभाजन चाहिए ---
अंडे से निकलने के लगभग तीन विभाजन बाद।
\textbf{9.} पॉलीमरेज़ लक्ष्य संदेशवाहक की नकल करके नया द्विलड़ी RNA
बनाती है, जो द्वितीयक \omterm{def:b3:rna-regulation:ncrna}{siRNA} में कटता है: उत्प्रेरक स्वयं लक्ष्य से फिर
बनता है, इसलिए वह तनुकरण से बचता और फैलता है (कोई वाहिका द्विलड़ी RNA
को कोशिकाओं के बीच और जनन-वंश में पहुँचाती है)। वह इसलिए मिटता है कि
प्रवर्धन के लिए लक्ष्य और कोई प्राथमिक उत्प्रेरक दोनों चाहिए; अंतःक्षेपित
RNA के जाते ही द्वितीयक भंडार पीढ़ी-दर-पीढ़ी तनु होता जाता है और
जनन-वंश शून्य हो जाता है।
\textbf{10.} $5000/2^{n} < 100$: $2^{n} > 50$, $n = 6$ दिन ($78$
प्रतियाँ बचती हैं)।
\textbf{11.} अविभाजनशील कोशिकाएँ \omterm{def:b3:rna-regulation:rnai}{RISC} को तनु नहीं करतीं; सीमा \omterm{def:b3:rna-regulation:ncrna}{siRNA} का
रासायनिक जीवन-काल है (न्यूक्लिएज़ से अपघटन, जिसे लड़ियों के रासायनिक
रूपांतरण से धीमा किया जाता है) और उस अंतःपुटीय भंडार का धीमा रिसाव जो
कोशिकाद्रव्य को भरता है।
\textbf{12.} $2\times 10^{7}/\num{16384} \approx \num{1200}$ संयोगी
मेल। निर्णय आनुवंशिकी ने किया: \emph{lin-14} का ह्रास \emph{lin-4} के
ह्रास से उलटा लक्षण-प्ररूप देता है, \emph{lin-14} के प्रकार्य-लाभ
युग्मविकल्पी ठीक 3$'$ UTR स्थलों के विलोपन थे, और वे सात स्थल संरक्षित थे।
\textbf{13.} स्तर $\propto$ अर्ध-आयु काल: $50/10 = 5$ गुना उठान।
\textbf{14.} फ़ेरिटिन संश्लेषण में $1/0.05 = 20$ गुना गिरावट।
\textbf{15.} लौह-समृद्ध से लौह-दरिद्र तक आयात/भंडारण $5\times 20 = 100$
गुना बढ़ता है।
\textbf{16.} गुच्छ तभी जुड़ता है जब कोशिकाद्रव्यी लोहा उपलब्ध हो,
इसलिए प्रोटीन की संरूपता लौह सांद्रता का सीधा पाठ है, बिना किसी ग्राही,
हॉर्मोन या अनुलेखन के --- किसी राइबोस्विच की तरह, ऐसा उपापचयज जिसे वही
अणु भाँपता है जिसे वह नियंत्रित करता है।
\textbf{17.} फ़ेरिटिन का अनुवादन लोहे की परवाह किए बिना होता है: सीरम
फ़ेरिटिन बहुत ऊँचा, सीरम लोहा और भंडार सामान्य (कोई अतिभार नहीं)। लेंस
में, जिसमें कोई आवागमन नहीं, फ़ेरिटिन वर्षों जमा होकर क्रिस्टलित हो
जाता है --- अगेता मोतियाबिंद।
\textbf{18.} हेयरपिन केवल उतरने की जगह है; प्रभाव वही है जो बँधा
प्रोटीन भौतिक रूप से रोकता है --- 5$'$ सिरे पर छानते राइबोसोमों को,
3$'$ सिरे पर किसी न्यूक्लिएज़ को।
\textbf{19.} दो रसायन दो स्थितियाँ ढकते हैं: IRP1 का लौह--गंधक गुच्छ
ऑक्सीकारकों और कम ऑक्सीजन से भी नष्ट होता है, इसलिए IRP1 अपचयोपचय
अवस्था पर अनुक्रिया करता है, जबकि IRP2 का लौह-निर्भर अपघटन शारीरिक
ऑक्सीजन परास भर में स्वयं लोहे पर अनुक्रिया करता है; बहुतायत
प्रकार्यगुच्छ को दृढ़ बनाती है, और दोनों भिन्न बातें बताते हैं।
\textbf{20.} $12\times 48\times 33\times 2 = \num{38016}$। $1 - (1 -
50/\num{38016})^{50} = 1 - 0.99868^{50} \approx 1 - e^{-0.066} =
0.064$: लगभग \qty{6}{\%}।
\textbf{21.} एक जैसे समरूपांतर समजातीय रूप से बँधकर प्रतिकर्षित करते
हैं: एक ही तंत्रिका-कोशिका की शाखाएँ एक-दूसरे से बचकर फैल जाती हैं,
जबकि भिन्न तंत्रिका-कोशिकाओं की शाखाएँ, जिनमें लगभग कोई समरूपांतर साझा
नहीं, आपस में कट सकती हैं --- प्रति कोशिका एक निजी पहचान।
\textbf{22.} $2^{n}$; $\num{1024}$, जब $n = 10$ हो। वास्तव में कम:
समावेश कोशिका-प्रकार-विशिष्ट कारक तय करते हैं और वह बहिर्वेशनों के पार
समन्वित होता है, अनेक संयोजन वाचन-ढाँचा खिसका देते हैं और
निरर्थकता-मध्यस्थ क्षय से नष्ट हो जाते हैं, और बहिर्वेशन स्वतंत्र नहीं हैं।
\textbf{23.} \emph{Sxl} \omterm{def:b3:chromatin-epigenetics:x-inactivation}{मात्रा प्रतिपूरण} भी नियंत्रित करता है: मादाओं
में SXL \emph{msl-2} का अनुवादन रोकता है, इसलिए दोनों X गुणसूत्रों का
अति-अनुलेखन नहीं होता; SXL के बिना कोई XX भ्रूण दोनों X गुणसूत्रों पर
MSL संकुल जोड़ लेता है और उनका उत्पादन दोगुना कर देता है --- हर
X-सहलग्न जीन की घातक अति-मात्रा, इससे पहले कि लिंग का प्रश्न भी उठे।
\textbf{24.} अगेता प्रवर्तक तभी तक चलता है जब तक X:A संकेत उपस्थित हो;
जो SXL वह बनाता है वह अनुरक्षण प्रवर्तक से आए अनुलेख का उत्पादक
विच्छेदन बाध्य कर देता है, और वह प्रवर्तक दोनों लिंगों में सक्रिय है,
इसलिए SXL एक बार बन जाए तो वह स्वयं को बनाता रहता है और संकेत की
आवश्यकता नहीं रहती --- वही पाठक-लेखक को बुलाता है वाला स्वयं-पोषी
तर्क जो \omterm{def:b3:chromatin-epigenetics:nucleosome}{क्रोमैटिन} चिह्नों का है, यहाँ एक प्रोटीन और एक विच्छेदन के साथ।
\textbf{25.} LIN-14 प्रोटीन $8$ गुना दमित; संदेशवाहक स्विच का
अर्ध-काल \qty{10}{min}; अप्रवर्धित कोई RNA $12$ विभाजनों के बाद प्रति
कोशिका एक अणु से नीचे गिर जाता है।
\end{solution}
